\documentclass[12pt]{article}
\usepackage[utf8]{inputenc}
\usepackage[T1]{fontenc}
\usepackage{hyperref}
\title{Pohnpeian Language Resources for K-12 Curriculum Development}
\author{Compiled by Claude Sonnet 5.5 (claude-sonnet-5-5) for Andr\'as Kornai}
\date{}
\begin{document}
\maketitle

\section{Introduction}
Pohnpeian is the main language of Pohnpei State, Federated States of Micronesia, written in a Latin orthography. Reference works and a few digital resources were found, but no official Department of Education curriculum, textbook list or standards for Pohnpeian appeared in search results.

\section{Primary Educational Resources}
\begin{itemize}
\item College of Micronesia-FSM course ML100, Introduction to Pohnpeian Language, developed mainly for in-service teachers to support Pohnpeian as a language of instruction and subject: \url{https://app.comfsm.fm/national/administration/VPIA/slo/ml/ML100\%20Introduction\%20to\%20Pohnpeian\%20Language.htm}
\item College of Micronesia-FSM course ML101P, Pohnpeian Reading: \url{https://app.comfsm.fm/national/administration/VPIA/slo/ml/ML101P\%20Pohnpeian\%20Reading.htm}
\end{itemize}

\section{Linguistic and Reference Materials}
\begin{itemize}
\item Rehg and Sohl, \emph{Ponapean-English Dictionary} (University of Hawaii Press, 1979), designed also for Ponapean educators: \url{https://uhpress.hawaii.edu/product/ponapean-english-dictionary/}
\item Rehg, \emph{Ponapean Reference Grammar} (1981), bibliographic record: \url{https://habeleinstitute.org/wiki/Zotero_reference:QSMLC8SE}
\item Wikipedia, Pohnpeian language (orthography and background): \url{https://en.wikipedia.org/wiki/Pohnpeian_language}
\item Omniglot, Pohnpeian writing: \url{https://www.omniglot.com/writing/pohnpeian.htm}
\end{itemize}

\section{Religious and Community Resources}
\begin{itemize}
\item Pohnpeian New Testament and Psalms (older orthography, per the portal): \url{https://dbs.org/bibles/PONPJA}
\end{itemize}

\section{Digital and Audio Resources}
\begin{itemize}
\item Pohnpeian Talking Dictionary (Living Tongues Institute, 2013; reuse requires permission): \url{https://talkingdictionary.swarthmore.edu/pohnpeian/}
\item Publication record for the talking dictionary: \url{https://works.swarthmore.edu/fac-linguistics/301}
\end{itemize}

\section{Curriculum Development Context}
Pohnpeian is a language of instruction in early grades, but no curriculum documents were located. Next steps would be to contact the Pohnpei State Department of Education and COM-FSM for classroom materials, and to use the dictionary and grammar above as references.
\end{document}